\chapter{Conclusioni}
Questo lavoro mi ha permesso di toccare con mano quali sono i rischi e le vulnerabilità ancora persistenti all'interno delle aziende, nonostante il grande aiuto fornito dall'IA nella generazione automatica di patch e rilevamento repentino delle minacce. Di vitale importanza per preservare i propri sistemi, in parallelo con la sensibilizzazione dell'utente sui possibili rischi, rimane fondamentale la preparazione dei team in incident response contro attacchi in corso o futuri ed il costante aggiornamento dell'infrastruttura.\\
Il framework Caldera\textsuperscript{TM}, grazie al suo costante sviluppo e aggiornamento da parte di MITRE, rimarrà uno strumento importante da utilizzare nel mercato dei cyber range anche nel prossimo decennio, permettendo ai componenti del Red Team di riuscire ad organizzare al meglio il lavoro, familiarizzandosi anche con la maggior parte dei threat reali, e di spostare la loro attenzione sull'ampliamento delle proprie conoscenze in modo da poter sostenere minacce che evolvono negli anni. È bene sottolineare che un insieme di tools come Caldera, ad esempio Mordor o Atomic Red Team, non riuscirà in nessun caso a sostituire completamente una squadra rossa con un bagaglio notevole di esperienza, questo poiché tale squadra sarà in grado di sfruttare vulnerabilità non ancora documentate in ATT\&CK. Un'intersezione nel ruolo da parte di entrambi, l'abbiamo nella fase di progettazione del cyber range: il Red Team evidenzia le aree di testing mirate a rappresentare una debolezza nella difesa della compagnia, mentre i framework possono essere configurati al fine di elevare il grado di autenticità con un attacco reale.\\[2em]
{\bfseries\Large{Possibili usi della VM}}\\[1.5em]
Una volta pubblicata, la VM Windows vulnerabile ha la possibilità di essere utilizzata in vari contesti e riconfigurata a seconda delle proprie esigenze:
\begin{description}
    \item[Host in un cyber range] $\rightarrow$ L'uso basico che si può fare della macchina, è lo stesso esplicitato in questa relazione. L'azienda che usufruisce della VM può riconfigurarla in modo da inserirla all'interno della propria rete ed includerla in qualche futuro cyber range in sviluppo. Ovviamente, è anche possibile migliorare o ampliare l'applicazione web installata in modo da renderla più realistica, aggiungendo magari dei servizi che la stessa enterprise offre ai clienti.
    \item[Challenge CTF] $\rightarrow$ Mi è stata data l'idea sulla creazione di una possibile sfida stile "Capture The Flag", in modo da rendere l'apprendimento di nozioni o il test delle proprie capacità ancora più interessanti. Può rappresentare una sfida personale, ma anche una possibile esercitazione da proporre ai propri dipendenti in modo da elevare l'attenzione verso i "rischi del mestiere". Molti sono i siti famosi in cui sono presenti queste challenge come HackTheBox o TryHackMe.
    \item[Honeypot] $\rightarrow$ Una debolezza come quella presentata all'interno dell'applicazione web diventa un'esca molto allettante per attirare dei possibili attaccanti ostili. L'uso di honeypots è molto diffuso soprattutto in grandi imprese e la suddetta macchina potrebbe entrare a far parte di questa feature, il tutto sempre configurandola secondo la policy dell'azienda cliente.
    \item[Test personali] $\rightarrow$ In base alla mia esperienza personale, ma soprattutto alle possibilità di ogni individuo, la continua esercitazione e il cimento nelle operazioni di cybersecurity è fondamentale al fine di acquisire dimestichezza con i tool più utilizzati. Ma ancora più importante, è riuscire a sviluppare dei progetti personali in questo campo che forniranno un grande trampolino di lancio verso lo sviluppo della propria carriera. Un uso consigliato dell'host è, quindi, quello di macchina per testare i software personali sviluppati o provare tecniche e tool appresi in un particolare corso.
    \item[Target esercitazioni con Caldera] $\rightarrow$ Lavorando con il framework, ho potuto esplorare diverse emulazioni di avversari contro un sistema vulnerabile; il numero delle configurazioni possibili con cui avviare un'operazione è notevolmente elevato. Non riuscendo a racchiuderle tutte in un unico documento, ritengo necessario offrire la VM Windows come strumento di "gioco" ed esercitazione con i plugins attualmente disponibili in Caldera; a questo scopo, si include un invito ad ampliare la propria infrastruttura di rete simulata, in modo da riuscire a usare le varie tecniche di \textit{lateral movement} non incluse durante le esercitazioni. 
    
\end{description}